A separately registered componentwise-delay policy accepts the already completed R1 plateau candidate on an independent branch, without changing the old strict-CPD rejection or rerunning native analysis. Its pure decision tests include acceptance and 14 invalid, nonmonotone, incomplete or repeated-state controls. The new policy requires a complete zero-RC physical audit, exactly one strictly shortened connection, all other sink delays unchanged within the audit tolerance, and nonworsening native setup/hold CPD/WNS/TNS at reported precision. This plateau step still has CPD 6.01810\,ns. Its next five current worst-path connections all have exact zero free-capacity improvement, so acceptance alone has not produced further CPD gain.

A same-target counterfactual isolates resource competition: all five target trees, roots, sinks, retained anchors, prefixes and physical delays are identical before and after the plateau step. Only six RR ownership entries change. Before the step, the next tied connection (net 6587 pin 4) has a 0.185826\,ns free-capacity improvement; afterwards it has none. Thus componentwise nonworsening present delays do not guarantee preservation of another connection's future routing opportunity. The physical shortening is real, but the first move reserves part of the corridor that the next tied connection would use.

A joint reference frees both target-exclusive chains while retaining all other branches and networks. Each target individually reaches 0.856978\,ns, but forbidding their shared unit-capacity IPIN 82120 raises each exact cost to its original 1.042804\,ns. Therefore both cannot simultaneously fall below that threshold with this fixed background. This conditional local minimax statement is not a global physical or STA bound. Removing all other occupancy and reducing each target tree to its SOURCE reveals 0.849339\,ns alternatives that avoid the selected resource; tests of other shared resources and one new conflict refinement produce three distinct paths per target. Of nine path pairs, two satisfy the two-route RR-node capacities, both with maximum cost 0.849339\,ns. This matches the independent single-endpoint lower bound and solves that optimistic two-endpoint minimax relaxation within numerical tolerance. It does not yet restore other terminals or networks.

Both selected paths also attach to the actual retained trees at unchanged cost and have no mutual capacity conflict. Their remaining conflicts involve four other networks and four affected sinks. A six-network transaction reserves the two target paths together and restores those four sinks in ascending net/pin order. It passes unchanged-native FULL legality, independent physical/capacity checks and full STA. Both target delays decrease from 1.042804 to 0.849339\,ns. Six sink delays change in total, with three improvements and three regressions; all unrelated network traces remain identical. Final CPD improves from the strong R1 value 6.01810 to 5.85466\,ns (2.715807\%). Setup TNS improves from $-33483.7$ to $-33438.8$\,ns, and hold WNS/TNS remain zero. The worst path moves to a clk-domain path from the FPAddSub AlignModule register to the FPMult ExecuteModule register.

Occupied SB pairs increase from 39 to 40, CB crossing output sources from 969 to 970, and planar wirelength from 218,164 to 218,176. The prototype constructs and physically audits the transaction in 37.576534\,s; four query cores sum to 1.344968\,s within four query processes totaling 4.819448\,s. The native load/check/STA process takes 8.945292\,s. These stage measurements are not formal complete-route timing or an integrated implementation cost. This study uses 31 new reference queries, four transaction queries and one new legal native analysis call; the original plateau native evidence is reused. The gain supports joint critical-group selection and coordinated resource restoration on this development case, not plateau acceptance alone or broad method superiority. Exact native all-connection timing/pin information and equal-opportunity complete-policy A0/R1/Tf comparisons are the next work.
